\chapterhead{1}{Moduli schemes and the q-expansion principle}

In this chapter, we will recall some of the definitions and main results of the
theory of moduli of elliptic curves\index{elliptic curve}, and deduce from them various forms of the
"q-expansion principle\index{q-expansion principle}" for modular forms\index{modular form}.

\kaitemb[1.0]{\uline{1.0.}} By an elliptic curve over a scheme $S$, we mean a proper smooth
morphism $p\colon E\to S$, whose geometric fibres are connected curves of genus
one, together with a section $e\colon S\to E$.

\begin{center}
\begin{tikzcd}
E \arrow[d, "p"'] \\
S \arrow[u, bend right, "e"']
\end{tikzcd}
\end{center}

We denote by $\w_{E/S}$ the invertible sheaf $p_*(\Omega^1_{E/S})$ on $S$,
which is canonically dual (Serre duality) to the invertible sheaf
$R^1p_*(\cO_E)$ on $S$.

\sechead{1.1}{Modular forms of level 1}

A modular form\index{modular form} of weight $k\in\Z$ and level one is a rule $f$ which assigns to
any elliptic curve\index{elliptic curve} $E$ over any scheme $S$ a section $f(E/S)$ of
$(\w_{E/S})^{\otimes k}$ over $S$ such that the following two conditions are
satisfied.

1. $f(E/S)$ depends only on the $S$-isomorphism class of the elliptic curve
$E/S$.

2. The formation of $f(E/S)$ commutes with arbitrary change of base
$g\colon S'\to S$ (meaning that $f(E_{S'}/S') = g^*f(E/S)$).

We denote by $M(\Z;1,k)$ the $\Z$-module of such forms.

Equivalently, a modular form of weight $k$ and level $1$ is a rule $f$ which
assigns to every pair $(E/R,\omega)$ consisting of an elliptic curve over (the
spectrum of) a ring $R$ together with a basis $\omega$ of $\w_{E/R}$ (i.e., a
nowhere vanishing section of $\Omega^1_{E/R}$ on $E$), an element
$f(E/R,\omega)\in R$, such that the following three conditions are satisfied.

1. $f(E/R,\omega)$ depends only on the $R$-isomorphism class of the pair
$(E/R,\omega)$.

2. $f$ is homogeneous of degree $-k$ in the "second variable"; for any
$\lambda\in R^\times$ (the multiplicative group of $R$),
$f(E,\lambda\omega) = \lambda^{-k}f(E,\omega)$.

3. The formation of $f(E/R,\omega)$ commutes with arbitrary extension of
scalars $g\colon R\to R'$ (meaning $f(E_{R'}/R',\omega_{R'}) =
g(f(E/R,\omega))$).

(The correspondence between the two notions is given by the formula
$f(E/\Spec(R)) = f(E/R,\omega)\cdot\omega^{\otimes k}$ valid whenever
$S=\Spec(R)$ and $\w_{E/R}$ is a \uline{free} $R$-module, with basis $\omega$.)

If, in the preceding definitions we consider \uline{only} schemes $S$ (or
rings $R$) lying over a \uline{fixed ground-ring} $R_0$, and only changes of
base by $R_0$-morphisms, we obtain the notion of a modular form of weight $k$
and level one defined over $R_0$, the $R_0$-module of which is noted
$M(R_0,1,k)$.

A modular form $f$ of weight $k$ and level one defined over $R_0$ can be
evaluated on the pair $(\mathrm{Tate}(q),\omega_{can})_{R_0}$ consisting of the
Tate curve\index{Tate curve} and its canonical differential, viewed as elliptic curve with
differential over $\Z((q))\tens_\Z R_0$ (and not just over $R_0((q))$).

The q-expansion of a modular form $f$ is by definition the finite-tailed
Laurent series

\[
f((\mathrm{Tate}(q),\omega_{can})_{R_0})\in\Z((q))\tens_\Z R_0.
\]

The modular form $f$ is called \uline{holomorphic at $\infty$} if its
q-expansion lies in the subring $\Z[[q]]\tens_\Z R_0$; the module of all such
is noted $S(R_0;1,k)$. Notice that the q-expansion lies in
$\Z((q))\tens_\Z R_0\subset R_0((q))$, i.e., it is \uline{finite}
$R_0$-linear combination of elements of $\Z((q))$. This implies, for example,
that if $R_0$ is the field of fractions of a discrete valuation ring, then the
q-expansion coefficients of any modular form of weight $k$ and level one over
$R_0$ have \uline{bounded denominators}.

\sechead{1.2}{Modular forms of level n}

For each integer $n\geq 1$, we denote by ${}_nE$ the kernel of "multiplication
by $n$" on $E/S$; it is a finite flat commutative group-scheme of rank $n^2$
over $S$, which is étale over $S$ if and only if the integer $n$ is invertible
in $\Gamma(S,\cO_S)$ i.e., if and only if $S$ is a scheme over
$\Z[\frac{1}{n}]$. A level $n$ structure on $E/S$ is an isomorphism

\[
\alpha_n\colon {}_nE\xrightarrow{\sim}(\Z/n\Z)^2_S.
\]

It cannot exist unless $n$ is invertible on $S$, and in that case there always
exists one on some finite étale covering $S'$ of $S$. If a level $n$ structure
on $E/S$ exists, and if $S$ is connected, the set of all such is principal
homogeneous under $GL(2,\Z/n\Z) = \Aut((\Z/n\Z)^2_S)$.

A modular form\index{modular form} of weight $k$ and level $n$ is rule which assigns to each pair
$(E/S,\alpha_n)$ consisting of an elliptic curve\index{elliptic curve} together with a level $n$
structure a section $f(E/S,\alpha_n)$ of $(\w_{E/S})^{\otimes k}$ over $S$, in
a way which depends only on the isomorphism class of $(E/S,\alpha_n)$, and
which commutes with arbitrary base-change $g\colon S'\to S$. Equivalently, it
is a rule which assigns to all triples $(E/R,\omega,\alpha_n)$, consisting of
an elliptic curve over a ring $R$ together with a base $\omega$ of $\w_{E/R}$
and a level $n$ structure $\alpha_n$, an element
$f(E/R,\omega,\alpha_n)\in R$ which depends only on the isomorphism class of
$(E/R,\omega,\alpha_n)$, which commutes with arbitrary change of base, and which
is homogeneous of degree $-k$ in the "second variable", meaning that for any
$\lambda\in R^\times$, we have $f(E/R,\lambda\omega,\alpha_n) =
\lambda^{-k}f(E/R,\omega,\alpha_n)$. Exactly as for level one, we define the
notion of a modular form of weight $k$ and level $n$ defined over a ring $R_0$.
The $R_0$-module of all such is noted $M(R_0,n,k)$.

A modular form of weight $k$ and level $n$ defined over a ring $R_0$ which
contains $1/n$ \uline{and} a primitive $n$'th root of unity $\zeta_n$ may be
evaluated on the triples $(\mathrm{Tate}(q^n),\omega_{can},\alpha_n)_{R_0}$
consisting of the Tate curve\index{Tate curve} $\mathrm{Tate}(q^n)$ with its canonical
differential, viewed as defined over $\Z((q))\tens_\Z R_0$, together with any
of its level $n$ structures (\uline{all} points of ${}_nE$ are rational over
$\Z((q))\tens_\Z R_0$; in fact, being the canonical images of the points
$\zeta_n^i q^j$, $0\leq i,j\leq n-1$ from "$\Gm$", they all have coordinates in
$\Z[[q]]\tens_\Z\Z[\frac{1}{n},\zeta_n]$, and the non-constant q-coefficients
of their $(x,y)$ coordinates even lie in $\Z[\zeta_n]$ (cf.\cite{roquette70}), as one sees
using the explicit formulas of Jacobi-Tate.

The q-expansions of the modular form $f$ are the finitely many finite-tailed
Laurent series

\begin{kadispb}{1.2.1}
f((\mathrm{Tate}(q^n),\omega_{can},\alpha_n)_{R_0})\in\Z((q))\tens_\Z R_0
\end{kadispb}

obtained by varying $\alpha_n$ over all the level $n$ structures.

(\underline{NB} Though it makes sense to speak of a modular form of weight $k$
and level $n$ defined over \uline{any} ring $R_0$, we can speak of its
q-expansions over $R_0$ \uline{only} when $R_0$ contains $1/n$ and a primitive
$n$'th root $\zeta_n$ of $1$.)

A modular form defined over any ring $R_0$ is said to be holomorphic at
$\infty$ if its inverse image on $R_0[1/n,\zeta_n]$ has \uline{all} its
q-expansions in $\Z[[q]]\tens_\Z R_0[1/n,\zeta_n]$. $\bigl\langle$ If the ring
$R_0$ itself contains $1/n$ and $\zeta_n$, this is equivalent to asking that
\uline{all} the q-expansions lie in $\Z[[q]]\tens_\Z R_0$. $\bigr\rangle$ The
module of such is denoted $S(R_0;n,k)$.

A modular form (resp: holo. at $\infty$) of weight $k$ and level $n$ defined
over a ring $R_0$, which \uline{does not depend} on the "last variable"
$\alpha_n$ \uline{is} a modular a modular form (resp: holo. at $\infty$) of
weight $k$ and level one defined over $R_0[1/n]$.

\sechead{1.3}{Modular forms on $\Gamma_0(p)$}

Analogously, for an integer $n\geq 1$ and a prime number $p\nmid n$, a modular
form of weight $k$ and level $n$ on $\Gamma_0(p)$ is a rule $f$ which assigns
to each triple $(E/S,\alpha_n,H)$ consisting of an elliptic curve\index{elliptic curve}, a level $n$
structure, and a finite flat subgroup-scheme $H\subset E$ of rank $p$, a
section $f(E/S,\alpha_n,H)$ of $(\w_{E/S})^{\otimes k}$ over $S$, which depends
only on the isomorphism class of $(E/S,\alpha_n,H)$, and whose formation
commutes with arbitrary change of base $S'\to S$. Equivalently, it is a rule
which assigns to each quadruple $(E/R,\omega,\alpha_n,H)$ an element
$f(E/R,\omega,\alpha_n,H)\in R$, which depends only on the isomorphism class of
the quadruple, whose formation commutes with arbitrary change of base, and
which is homogeneous of degree $-k$ in the second variable. As before, we
define the notion of a modular form\index{modular form} of weight $k$ and level $n$ on
$\Gamma_0(p)$ being defined over a ring $R_0$.

A modular form of weight $k$ and level $n$ on $\Gamma_0(p)$, defined over a ring
$R_0$ which contains $1/n$ and $\zeta_n$ may be evaluated on each of the
quadruples $(\mathrm{Tate}(q^n),\omega_{can},\alpha_n,\mu_p)_{R_0}$. We will
call the values of $f$ on these quadruples the \uline{q-expansions} of $f$ at
the unramified cusps, and say that $f$ is holomorphic at the unramified cusps
if its q-expansions there all lie in $\Z[[q]]\tens_\Z R_0$. We can also
evaluate $f$ on each of quadruples
$(\mathrm{Tate}(q^{np}),\omega_{can},\alpha_n,\{q^n\})$, where $\{q^n\}$
denotes the flat rank-p subgroup scheme generated by (the image of) $q^n$. Its
values there are called its q-expansions at the ramified cusps. We say that $f$
is holomorphic at $\infty$ if \uline{all} of its q-expansions, at the ramified
and unramified cusps, actually lie in $\Z[[q]]\tens_\Z R_0$.

\begin{kathmstar}{Remark.}
The distinction between ramified and unramified cusps on $\Gamma_0(p)$ is
quite a natural one - in the work of Atkin, one deals with modular functions
(weight 0) of level one on $\Gamma_0(p)$ which are holomorphic at the
unramified cusp, but \uline{not} at the ramified one.
\end{kathmstar}

\sechead{1.4}{The modular schemes $M_n$ and $\Mbar_n$}

For each integer $n\geq 3$, the functor "isomorphism classes of elliptic curves\index{elliptic curve}
with level n structure" is \uline{representable}, by a scheme $M_n$ which is an
affine smooth curve over $\Z[\frac{1}{n}]$, finite and flat of degree
$=\#(\GL_2(\Z/n\Z)/\pm 1)$ over the affine j-line $\Z[\frac{1}{n},j]$, and
étale over the open set of the affine j-line where $j$ and $j-1728$ are
invertible. The \uline{normalization} of the projective j-line
$\PP^1_{\Z[1/n]}$ in $M_n$ is a proper and smooth curve $\Mbar_n$ over
$\Z[1/n]$, the global sections of whose structural sheaf are
$\Z[1/n,\zeta_n]$. The curve $M_n\tens_{\Z[1/n]}\Z[1/n,\zeta_n]$ (resp.
$\Mbar_n\tens_{\Z[1/n]}\Z[1/n,\zeta_n]$) is a disjoint union of $\varphi(n)$
affine (resp. proper) smooth geometrically connected curves over
$\Z[1/n,\zeta_n]$, the partitioning into components given by the $\varphi(n)$
primitive $n$'th roots of one occurring as values of the e.m. pairing on the
basis of ${}_nE$ specified by the level n structure. The scheme $\Mbar_n - M_n$
over $\Z[1/n]$ is finite and étale, and over $\Z[1/n,\zeta_n]$, it is a
disjoint union of sections, called the \uline{cusps} of $\Mbar_n$, which in a
natural way \uline{are} the set of isomorphism classes of level n structures on
the Tate curve\index{Tate curve} $\mathrm{Tate}(q^n)$ viewed over
$\Z((q))\tens_\Z\Z[1/n,\zeta_n]$. The completion of
$\Mbar_n\tens\Z[1/n,\zeta_n]$ along any of the cusps is isomorphic to
$\Z[1/n,\zeta_n][[q]]$. The completion of the projective j-line
$\PP^1_{\Z[1/n,\zeta_n]}$ along $\infty$ is itself isomorphic to
$\Z[1/n,\zeta_n][[q]]$, via the formula
$j(\mathrm{Tate}(q)) = 1/q + 744 + \ldots$, and the endomorphism of
$\Z[1/n,\zeta_n][[q]]$ arising from the projection $\Mbar_n\to\PP^1$ is just
given by $q\mapsto q^n$. In fact, for each cusp, the inverse image of the
\uline{universal} elliptic curve with level n structure $(E/M_n,\alpha_n)$ over
(the spectrum of) $\Z[1/n,\zeta_n]((q))$ (viewed as a punctured disc around the
cusp) is isomorphic to the inverse image over $\Z[1/n,\zeta_n]((q))$ of the
Tate curve $\mathrm{Tate}(q^n)$ with the level n structure corresponding to
that cusp.

\sechead{1.5}{The invertible sheaf $\w$ on $\Mbar_n$, and modular forms holomorphic at $\infty$}

There is a unique invertible sheaf $\w$ on $\Mbar_n$ whose restriction to $M_n$
is $\w_{E/M_n}$ ($(E/M_n,\alpha_n)$ the universal elliptic curve\index{elliptic curve} with level n
structure), and whose sections over the completion $\Z[1/n,\zeta_n][[q]]$ at
each cusp are precisely the $\Z[1/n,\zeta_n][[q]]$ multiples of the canonical
differential of the Tate curve\index{Tate curve}. The Kodaira-Spencer style isomorphism (cf.
\ref{ka:A1.3.17} and \cite{deligne69})

\[
(\w_{E/M_n})^{\otimes 2}\xrightarrow{\sim}\Omega^1_{M_n/\Z[1/n]}
\]

\uline{extends} to an isomorphism

\[
(\w)^{\otimes 2}\xrightarrow{\sim}\Omega^1_{\Mbar_n/\Z[1/n]}(\log(\Mbar_n-M_n)),
\]

and, in fact, over $\Z[1/n,\zeta_n][[q]]$, the "square" of the canonical
differential $\omega_{can}$ on $\mathrm{Tate}(q^n)$ corresponds to
$n\cdot\frac{\dd q}{q}$.

It follows that a modular form\index{modular form} of level $n$ and weight $k$ holomorphic at
$\infty$ defined over any ring $R_0\ni 1/n$ is just a section of
$(\w)^{\otimes k}$ on $\Mbar_n\tens_{\Z[1/n]}R_0$, or equivalently a section of
the quasi-coherent sheaf $(\w)^{\otimes k}\tens_{\Z[1/n]}R_0$ on $\Mbar_n$.

\sechead{1.6}{The q-expansion principle}

For \uline{any} $\Z[1/n]$-module $K$, we define a modular form\index{modular form} of level $n$ and
weight $k$, holomorphic at $\infty$, with coefficients in $K$, to be an element
of $H^0(\Mbar_n,(\w)^{\otimes k}\tens_{\Z[1/n]}K)$. At each cusp, such a
modular form has a q-expansion in
$K\tens_{\Z[1/n]}\Z[1/n,\zeta_n]\tens_\Z\Z[[q]]$.

\begin{kathm}{Theorem}{1.6.1.}
Let $n\geq 3$, $K$ a $\Z[1/n]$-module, and $f$ a modular form of level $n$ and
weight $k$, holomorphic at $\infty$, with coefficients in $K$. Suppose that on
\uline{each} of the $\varphi(n)$ connected components of
$\Mbar_n\tens_{\Z[1/n]}\Z[1/n,\zeta_n]$, there is at least \uline{one} cusp at
which the q-expansion of $f$ vanishes identically. Then $f = 0$.
\end{kathm}

Before proving it, we give the main corollary.

\begin{kathm}{Corollary}{1.6.2.}
(The q-expansion principle\index{q-expansion principle}). Let $n\geq 3$, $K$ a $\Z[1/n]$-module, $L\subset K$
a $\Z[1/n]$-submodule. Let $f$ be a modular form of weight $k$, level $n$,
holomorphic at $\infty$, with coefficients in $K$. Suppose that on \uline{each}
of the $\varphi(n)$ connected components of
$\Mbar_n\tens_{\Z[1/n]}\Z[1/n,\zeta_n]$, there is at least one cusp at which all
the q-coefficients of $f$ lie in $L\tens_{\Z[1/n]}\Z[1/n,\zeta_n]$. Then $f$
\uline{is} a modular form with coefficients in $L$.
\end{kathm}

\begin{kaproof}[Proof of corollary]
The exact sequence $0\to L\to K\to K/L\to 0$ of $\Z[1/n]$-modules gives an
exact sequence of sheaves on $\Mbar_n$,

\begin{kadispb}{1.6.2.1}
0\to L\tens(\w)^{\otimes k}\to K\tens(\w)^{\otimes k}\to (K/L)\tens(\w)^{\otimes k}\to 0,
\end{kadispb}

hence a cohomology exact sequence

\begin{kadispb}{1.6.2.2}
0\to H^0(\Mbar_n,L\tens\w^{\otimes k})\to H^0(\Mbar_n,K\tens(\w)^{\otimes k})\to H^0(\Mbar_n,(K/L)\tens(\w)^{\otimes k}).
\end{kadispb}

The theorem \eqref{ka:1.6.1} now applies to the \uline{image} of $f$ in
$H^0(\Mbar_n,(K/L)\tens\w^{\otimes k})$, showing that image to be zero, whence
$f\in H^0(\Mbar_n,L\tens(\w)^{\otimes k})$ by the cohomology exact sequence.
\end{kaproof}

We now turn to the proof of the theorem. By considering the \uline{ring} of
dual numbers on $K$, $D(K) = \Z[1/n]\oplus K$, [multiplication $(a,k)(a',k') =
(aa',ak'+a'k)$] we are reduced to the case where $K$ is a \uline{ring} over
$\Z[1/n]$. Because the formation of the cohomology of quasi-coherent sheaves on
quasi-compact schemes commutes with inductive limits, we are first reduced to
the case where $K$ is a \uline{finitely generated} ring over $\Z[1/n]$, then to
the case when $K$ is a noetherian local ring. By faithful flatness of the
completion, we further reduce to the case when $K$ is a \uline{complete}
Noetherian local ring, then by Grothendieck's comparison theorem to the case
when $K$ is an artin local ring. By Krull's intersection theorem, $f$ induces
the zero-section of $(\w)^{\otimes k}$ over an open neighborhood of at least one
cusp on each connected component of
$\Mbar_n\tens K\tens\Z[1/n,\zeta_n]$, hence on an \uline{open dense set} in
$\Mbar_n\tens K$. If $f$ is not zero, there exists a non-void closed subset $Z$
of $\Mbar_n\tens K$, containing no maximal point of $\Mbar_n\tens K$, on which
$f$ is supported. Over the local ring in $\Mbar_n\tens K$ of any maximal point
$z$ of $Z$, $f$ becomes non-canonically a section of
$\hat\cO_{z,\Mbar_n\tens K}$ which is supported at the closed point, i.e. for
any element $g\in\mathfrak{m}_z$ (the maximal ideal of $\hat\cO_{z,\Mbar_n\tens K}$),
there exists a power $g^n$ of $g$ such that $g^nf = 0$. Thus \uline{every}
element of $\mathfrak{m}_z$ is a zero-divisor, i.e. the point $z\in\Mbar_n\tens K$ has
\uline{depth zero}. As $\Mbar_n\tens K$ is \uline{smooth} over an artin local
ring $K$, it is Cohen-Macaulay, and hence only its maximal points have depth
zero. Thus $z$ must be a maximal point of $\Mbar_n\tens K$, a contradiction.
Hence $f$ must be zero. \hspace*{\fill}$\square$

\sechead{1.7}{Base-change of modular forms of level $n\geq 3$}

\begin{kathm}{Theorem}{1.7.1.}
Let $n\geq 3$, and suppose either that $k\geq 2$ \uline{or} that $k=1$ and
$n\leq 11$. Then for any $\Z[1/n]$-module $K$, the canonical map

\[
K\tens H^0(\Mbar_n,(\w)^{\otimes k})\to H^0(\Mbar_n,K\tens(\w)^{\otimes k})
\]

is an isomorphism.
\end{kathm}

\begin{kaproof}
By standard base-changing theorems, it suffices to show that
$H^1(\Mbar_n,\w^{\otimes k}) = 0$. The isomorphism
$(\w)^{\otimes 2}\xrightarrow{\sim}\Omega^1_{\Mbar_n/\Z[1/n]}(\log(\Mbar_n-M_n))$,
together with the fact that each connected component of
$\Mbar_n\tens\Z[1/n,\zeta_n]$ contains at least one cusp, shows that for
$k\geq 2$, the restriction of $(\w)^{\otimes k}$ to each connected component of
$\Mbar_n\tens\Z[1/n,\zeta_n]$ has \uline{degree strictly greater} than $2g-2$,
$g$ the (common) genus of any of these components, and hence
$H^1(\Mbar_n,(\w)^{\otimes k})=0$ by Riemann-Roch. For $3\leq n\leq 11$,
explicit calculation shows that $\w$ restricted to each connected component of
$\Mbar_n\tens\Z[1/n,\zeta_n]$ has degree strictly greater than $2g-2$, and we
conclude as before.
\end{kaproof}

\begin{kathmstar}{Remark.}
When $n\geq 12$, $\w$ has degree $\leq 2g-2$ on each connected component of
$\Mbar_n\tens\Z[1/n,\zeta_n]$, and equality holds \uline{only} for $n=12$. The
author does not know whether or not the formation of modular forms\index{modular form} of weight
\uline{one} and level $n\geq 12$ commutes with base change.
\end{kathmstar}

\sechead{1.8}{Base change of modular forms of level 1 and 2}

\begin{kathm}{Theorem}{1.8.1.}
Let $R_0$ be any ring in which $2$ is invertible. For every integer $k\geq 1$,
the canonical map $S(\Z,2,k)\tens_\Z R_0\to S(R_0,2,k)$ is an isomorphism.
\end{kathm}

\begin{kaproof}
First we should remark that there are no non-zero modular forms\index{modular form} of level two and
\uline{odd} weight $k$ over $R_0$, because the automorphism "$-1$" of an
elliptic curve\index{elliptic curve} transforms $(E,\omega,\alpha_2)$ into $(E,-\omega,-\alpha_2)$,
hence $f(E,\omega,\alpha_2) = f(E,-\omega,-\alpha_2)$, but $\alpha_2 = -\alpha_2$,
hence $f(E,-\omega,-\alpha_2) = f(E,-\omega,\alpha_2) = (-1)^{-k}f(E,\omega,\alpha_2)$,
hence $2f(E,\omega,\alpha_2) = 0$ for $k$ odd.

In any case, modular forms of level two and weight $k$, holomorphic at
infinity, over any ring $R_0\ni 1/2$, are precisely those modular forms of
level \uline{four} and weight $k$ holomorphic at $\infty$, defined over $R_0$,
which are \uline{invariant} under the action of the subgroup of
$\GL_2(\Z/4\Z)$ consisting of the matrices $\equiv I\bmod 2$. As this group has
order $16$, a power of two, we may simply apply the \uline{projector}
$\frac{1}{16}\sum_{g\equiv I(2)}g$ to the base-changing isomorphism \eqref{ka:1.7.1} in
level four to produce the desired isomorphism in level two.
\end{kaproof}

\begin{kathm}{Theorem}{1.8.2.}
Let $R_0$ be any ring in which $2$ \uline{and} $3$ are invertible. For every
integer $k\geq 1$, the canonical map

\[
S(\Z,1,k)\tens_\Z R_0\to S(R_0,1,k)
\]

is an isomorphism.
\end{kathm}

\begin{kaproof}
The proof is similar to the previous one. We view a modular form of level one
over a ring $R_0\ni 1/6$ as a modular form of level four (resp. three) invariant
under $\GL(2,\Z/4\Z)$ (resp. $\GL(2,\Z/3\Z)$), defined over $R_0$. As,
$\GL(2,\Z/4\Z)$ has order $96 = 32\times 3$ (resp. $\GL(2,\Z/3\Z)$ has order
$48 = 16\times 3$), the projection technique \eqref{ka:1.8.1} shows that the canonical
map

\[
S(\Z[1/6],1,k)\tens_{\Z[1/6]}R_0\to S(R_0,1,k)
\]

is an isomorphism. Thus it remains only to handle the passage from
$\Z[1/6]$. But for \uline{any} ring $R$, $S(R,1,k)$ is the fibre product of the
diagram:

\begin{equation}\tag{1.8.2.1}\label{ka:1.8.2.1}
\begin{tikzcd}
H^0(\Mbar_3\tens R,(\w)^{\otimes k}) \arrow[d] \\
H^0(\Mbar_{12}\tens R,(\w)^{\otimes k}) & H^0(\Mbar_4\tens R,(\w)^{\otimes k}) \arrow[l]
\end{tikzcd}
\end{equation}

(i.e. a modular form of level one over $R$ \uline{is} a modular from $f_3$ of
level three over $R[1/3]$ together with a modular form $f_4$ of level four over
$R[1/2]$, such that $f_3$ and $f_4$ induce the \uline{same} modular form of
level $12$ over $R[1/12]$). As the formation of the diagram \eqref{ka:1.8.2.1}
\uline{and} of its fibre product commutes with any flat extension of scalars
$R\to R'$, taking $R=\Z$, $R'=\Z[1/6]$ gives the desired result.
\end{kaproof}

\begin{kathmstar}{Remark 1.8.2.2.}
The above theorem becomes \uline{false} when we do not exclude the primes $2$
and $3$. For over the finite field $\F_p$, the Hasse invariant\index{Hasse invariant} $A$ \uline{is} a
modular form of level one and weight $p-1$, holomorphic at $\infty$. But over
$\Z$ there are \uline{no} non-zero modular forms over $\Z$ of level one,
holomorphic at $\infty$, of weight either one or two. Similarly, $A\cdot\Delta$
is a cusp form of level one and weight $13$ (resp. $14$) over $\F_2$ (resp.
$\F_3$), which cannot be the reduction mod p of a modular form over $\Z$. See
\cite{deligne68} for the full determination of modular forms over $\Z$.
\end{kathmstar}

\sechead{1.9}{Modular forms of level 1 and 2: q-expansion principle}

For $n=1$, $2$, and any $\Z[1/n]$-module $K$, we \uline{define} a modular form\index{modular form}
of level $n$ and weight $k$, holomorphic at $\infty$, with coefficients in $K$
to be

for $n=1$: an element of the fibre-product of the diagram

\begin{equation}\tag{1.9.0.0}\label{ka:1.9.0.0}
\resizebox{\textwidth}{!}{%
\begin{tikzcd}[ampersand replacement=\&]
H^0(\Mbar_3,(\w)^{\otimes k}\tens_{\Z[1/3]}(K\tens_\Z\Z[1/3])) \arrow[d] \\
H^0(\Mbar_{12},(\w)^{\otimes k}\tens_{\Z[1/12]}(K\tens_{\Z[1/12]}\Z[1/12])) \& H^0(\Mbar_4,(\w)^{\otimes k}\tens_{\Z[1/4]}(K\tens\Z[1/4])) \arrow[l]
\end{tikzcd}}
\end{equation}

\kaitem{1.9.0.1} for $n=2$: an element of
$H^0(\Mbar_4,(\w)^{\otimes k}\tens_{\Z[1/4]}K)$ invariant by the subgroup of
$\GL_2(\Z/4\Z)$ consisting of matrices $\equiv I\bmod 2$.

The module of all such is noted $S(K,n,k)$.

(In the case $K$ is a \uline{ring}, this notion coincides with that already
introduced.) An exact sequence $0\to L\to K\to K/L\to 0$ gives an exact sequence
(\uline{without} the final $0$) of modules of modular forms, analogous to
\eqref{ka:1.6.2.2}.

As a corollary of \eqref{ka:1.6.1}, we have

\begin{kathm}{Corollary}{1.9.1.}
(q-expansion principle\index{q-expansion principle}) Let $n=1$ or $2$, $K$ a $\Z[1/n]$-module, and
$L\subset K$ a $\Z[1/n]$ submodule. Let $f$ be a modular form of weight $k$,
level $n$, holomorphic at $\infty$, with coefficients in $K$. Suppose that at
\uline{one} of the cusps (for $n=1$, there is only one, $j=\infty$, while for
$n=2$ there are three, $\lambda = 0,1,\infty$), the q-coefficients of $f$ all
lie in $L$. Then $f$ \uline{is} a modular form with coefficients in $L$.
\end{kathm}

\sechead{1.10}{Modular schemes of level 1 and 2}

They don't exist, in the sense that the corresponding functors are not
representable. However, for each $n\geq 3$ we can form the quotients

\[
M_n/\GL_2(\Z/n\Z) = \text{the affine } j\text{-line } A^1_{\Z[1/n]}
\]
\[
\Mbar_n/\GL_2(\Z/n\Z) = \text{the projective } j\text{-line } \PP^1_{\Z[1/n]}
\]

which fit together for variable $n$ to form the affine and projective j-lines
over $\Z$. We \uline{define} $M_1 = A^1_\Z$, the affine j-line, and
$\Mbar_1 = \PP^1_\Z$. The invertible sheaf $\w$ on $\Mbar_n$, $n\geq 3$, does
\uline{not "descend"} to an invertible sheaf on $\Mbar_1$, but its $12^{th}$
power $\w^{\otimes 12}$ \uline{does} descend, to $\cO(1)$, the inverse of the
ideal sheaf of $\infty$.

In particular, modular forms\index{modular form} over any ring $R$ of level one and weight $12\cdot k$
holomorphic at $\infty$, are just the elements of $H^0(\PP^1_R,\cO(k))$, and
\uline{their} formation does commute with arbitrary change of base.

Analogously for $n=2$, we define

\[
M_2 = M_4/\text{the subgroup of } \GL_2(\Z/4\Z) \text{ of matrices } \equiv I\bmod 2
\]
\[
\Mbar_2 = \Mbar_4/\text{the subgroup of } \GL_2(\Z/4\Z) \text{ of matrices } \equiv I\bmod 2.
\]

The scheme $M_2$ is $\Spec\Z[\lambda][1/2\lambda(1-\lambda)]$, and $\Mbar_2$ is
the projective $\lambda$-line $\PP^1_{\Z[1/2]}$. The invertible sheaf $\w$ does
not descend to $\Mbar_2$, but its \uline{square} does descend, to $\cO(1)$ =
the inverse of the ideal sheaf of the cusp $\lambda = \infty$.

In particular, modular forms of level two over any ring $R\ni 1/2$, of
(necessarily!) \uline{even} weight $2k$ and holomorphic at all three cusps,
are just the elements of $H^0(\PP^1_R,\cO(k))$; hence their formation commutes
with arbitrary change of base.

\sechead{1.11}{Hecke operators}

Let $\ell$ be a prime number, $R$ a ring in which $\ell$ is invertible, and $n$
an integer prime to $\ell$. For any elliptic curve\index{elliptic curve} $E/R$, the group-scheme
${}_\ell E$ of points of order $\ell$ is finite étale over $R$, and on a finite
étale over-ring $R'$ it becomes non-canonically isomorphic to $(\Z/\ell\Z)^2_R$.
Thus over $R'$, the elliptic curve $E_{R'}/R'$ has precisely $\ell+1$ finite
flat subgroups-(schemes) of rank $\ell$. For any such subgroup $H$, we denote
by $\pi\colon E_{R'}\to E_{R'}/H$ the projection onto the quotient and by
$\check\pi\colon E_{R'}/H\to E_{R'}$ the \uline{dual} map, which is also finite
étale of degree $\ell$. The composition $\pi\cdot\check\pi$ is multiplication by
$\ell$ on $E_{R'}/H$, and the composition $\check\pi\cdot\pi$ is multiplication
by $\ell$ on $E_{R'}$.

If $\omega$ is a nowhere vanishing differential on $E/R$, then
$\check\pi^*(\omega_{R'}) = \mathrm{trace}_{\check\pi}(\omega_{R'})$ is a nowhere
vanishing differential on $E_{R'}/H$. If
$\alpha_n\colon {}_nE\xrightarrow{\sim}(\Z/n\Z)^2_R$ is a level $n$ structure on
$E/R$, there is unique level $n$ structure $\pi(\alpha_n)$ on $E_{R'}/H$ such
that the diagram

\begin{equation}\tag{1.11.0.0}\label{ka:1.11.0.0}
\begin{tikzcd}
& (\Z/n\Z)_{R'} \\
{}_nE_{R'} \arrow[ur, "\alpha_n"] \arrow[rr, "\pi"] & & {}_n(E_{R'}/H) \arrow[ul, "\pi(\alpha_n)"']
\end{tikzcd}
\end{equation}

is commutative. (N.B. There is another "natural" choice of level n structure on
$E_{R'}/H$, namely $\alpha_n\circ\check\pi = \ell\cdot\pi(\alpha_n)$, which we
will not use.)

Given a modular form\index{modular form} over $R$ of level $n$ and weight $k$, for each triple
$(E/R,\omega,\alpha_n)$ we may form the sum over the $\ell+1$ subgroups $H$ of
order $\ell+1$ of $E_{R'}$,

\begin{kadisp}{1.11.0.1}
\sum_H f(E_{R'}/H,\check\pi^*(\omega),\pi(\alpha_n))
\end{kadisp}

which, while apparently an element of $R'$, is in fact an element of $R$, and
does not depend on the auxiliary choice of $R'$. Normalizing this sum by the
factor $\ell^{k-1}$, we define the Hecke operator\index{Hecke operator} $T_\ell$ on modular forms of
level $n$ and weight $k$ by the formula

\begin{kadisp}{1.11.0.2}
(T_\ell f)(E/R,\omega,\alpha_n) = \ell^{k-1}\sum f(E_{R'}/H,\check\pi^*(\omega),\pi(\alpha_n)),
\end{kadisp}

the sum extended to the $\ell+1$ subgroups of order $\ell$.

We now consider the effect on the q-expansions. The $\ell$-division points of
the Tate curve\index{Tate curve} $\mathrm{Tate}(q^n)$ over $\Z((q))\tens_\Z\Z[1/n\ell]$ all become
rational over $\Z((q^{1/\ell}))\tens_\Z\Z[1/n\ell,\zeta_\ell]$, and the $\ell+1$
subgroups of order $\ell$ are the following:

\[
\mu_\ell, \text{ generated by } \zeta_\ell
\]
\[
H_i, \text{ generated by } (\zeta_\ell^i q^{1/\ell})^n \text{ for } i=0,1,\ldots,\ell-1.
\]

For the subgroup $\mu_\ell$, the quotient $\mathrm{Tate}(q^n)/\mu_\ell$ is
$\mathrm{Tate}(q^{n\ell})$ (the projection induced by the $\ell$'th power map on
$\Gm$) and the \uline{dual} isogeny consists of dividing $\mathrm{Tate}(q^{n\ell})$
by the subgroup generated by $q^n$. For the subgroups $H_i$, the quotient
$\mathrm{Tate}(q^n)/H_i$ is $\mathrm{Tate}((\zeta_\ell^i q^{1/\ell})^n)$, and the
dual isogeny consists of dividing $(\mathrm{Tate}((\zeta_\ell^i q^{1/\ell})^n)$
by \uline{its} subgroup $\mu_\ell$.

Thus for the subgroup $\mu_\ell$, we have
$\check\pi^*(\omega_{can}) = \omega_{can}$ on $\mathrm{Tate}(q^{n\ell})$, while
for the subgroups $H_i$, $\check\pi^*(\omega_{can}) = \ell\cdot(\omega_{can})$
on $\mathrm{Tate}((\zeta_\ell^i q^{1/\ell})^n)$ (because in the latter case
$\check\pi$ is induced by the $\ell$'th power mapping on $\Gm$, on which
$\omega_{can}$ is $\dd t/t$).

The quotient $\mathrm{Tate}(q^n)/\mu_\ell\xrightarrow{\sim}\mathrm{Tate}(q^{n\ell})$
may be viewed as obtained from $\mathrm{Tate}(q^n)$ by the extension of scalars
$\varphi_\ell\colon\Z((q))\to\Z((q))$ sending $q\mapsto q^\ell$. We denote by
$\alpha_n'$ the unique level n structure on $\mathrm{Tate}(q^n)$ such that
$\varphi_\ell^*(\alpha_n') = \pi_\ell(\alpha_n)$, $\pi_\ell(\alpha_n)$ denoting
the image of $\alpha_n$ by the \uline{projection} of $\mathrm{Tate}(q^n)$ onto
$\mathrm{Tate}(q^n)/\mu_\ell\xrightarrow{\sim}\mathrm{Tate}(q^{n\ell})$.

The quotients $\mathrm{Tate}(q^n)/H_i\xrightarrow{\sim}\mathrm{Tate}(q^{n/\ell}\zeta_\ell^{ni})$,
$i=0,\ldots,\ell-1$ over $\Z[1/n\ell,\zeta_{n\ell}]((q^{1/\ell}))$, may each be
viewed as obtained from $\mathrm{Tate}(q^n)/H_0\xrightarrow{\sim}\mathrm{Tate}(q^{n/\ell})$
by the extension of scalars
$\varphi_i\colon\Z[1/n\ell,\zeta_{n\ell}]((q^{1/\ell}))\allowbreak\to\Z[1/n\ell,\zeta_{n\ell}]((q^{1/\ell}))$
which sends $q^{1/\ell}\mapsto\zeta_\ell^i q^{1/\ell}$. Under this
identification, we have (noting $\pi_i\colon\mathrm{Tate}(q^n)\to\mathrm{Tate}(q^n)/H_i$,
$i=0,\ldots,\ell-1$ the projections) the relation
$\pi_i(\alpha_n) = \varphi_i^*(\pi_0(\alpha_n))$, as an immediate explicit
calculation shows. We denote by $\alpha_n''$ the level n structure
$i_\ell^*(\pi_0(\alpha_n))$ on $\mathrm{Tate}(q^n)$ obtained from
$\pi_0(\alpha_n)$ on $\mathrm{Tate}(q^{n/\ell})$ by the extension of scalars
$i_\ell\colon\Z[1/n\ell,\zeta_{n\ell}]((q^{1/\ell}))\xrightarrow{\sim}\Z[1/n\ell,\zeta_{n\ell}]((q))$
sending $q^{1/\ell}$ to $q$.

Thus we have

\begin{kadisp}{1.11.0.3}
f(\mathrm{Tate}(q^n)/\mu_\ell,\check\pi_\ell^*(\omega_{can}),\pi_\ell(\alpha_n)) &= f(\mathrm{Tate}(q^{n\ell}),\omega_{can},\varphi_\ell^*(\alpha_n')) \\
&= \varphi_\ell(f(\mathrm{Tate}(q^n),\omega_{can},\alpha_n')).
\end{kadisp}

\begin{kadisp}{1.11.0.4}
f(\mathrm{Tate}(q^n)/H_i,\check\pi_i^*(\omega_{can}),\pi_i(\alpha_n)) &= f(\mathrm{Tate}((\zeta_\ell^i q^{1/\ell})^n),\ell\cdot\omega_{can},\varphi_i^*(\pi_0(\alpha_n))) \\
&= \varphi_i(f(\mathrm{Tate}(q^{n/\ell}),\ell\cdot\omega_{can},\pi_0(\alpha_n))) \\
&= \varphi_i\circ(i_\ell)^{-1}(f(\mathrm{Tate}(q^n),\ell\cdot\omega_{can},\alpha_n'')) \\
&= \ell^{-1}\cdot\varphi_i\circ(i_\ell)^{-1}(f(\mathrm{Tate}(q^n),\omega_{can},\alpha_n'')).
\end{kadisp}

Combining these, we have the following formula for $T_\ell$.

\begin{kathm}{Formula}{1.11.1.}
Let $f$ be a modular form of level $n$ and weight $k$ over a ring $R$, and
suppose $\ell$ is a prime number not dividing $n$ which is invertible in $R$.
Let $f$ be a modular form of level $n$ and weight $k$, with q-expansions

\begin{kadisp}{1.11.1.0}
f(\mathrm{Tate}(q^n),\omega_{can},\alpha_n) = \sum_{i>-\infty}a_i(\alpha_n)\cdot q^i.
\end{kadisp}

Then

\begin{kadisp}{1.11.1.1}
(T_\ell f)(\mathrm{Tate}(q^n),\omega_{can},\alpha_n) = \sum_{i>-\infty}b_i(\alpha_n)q^i,
\end{kadisp}

where the coefficients $b_i(\alpha_n)$ are given by the formula

\begin{kadisp}{1.11.1.2}
b_i(\alpha_n) = \ell^{k-1}a_{i/\ell}(\alpha_n') + a_{\ell i}(\alpha_n'')
\end{kadisp}

(with the convention that $a_{i/\ell} = 0$ unless $\ell\mid i$).
\end{kathm}

\begin{kathm}{Corollary}{1.11.2.}
If $f$ is holomorphic at $\infty$, so is $T_\ell(f)$. If $f$ is a cusp-form
(meaning that its q-expansions all start in degree $\geq 1$), then so is
$T_\ell(f)$. If all the q-expansions of $f$ are \uline{polynomials} in $q$, the
same is true of $T_\ell(f)$.
\end{kathm}

\begin{kaproof}
These follow directly from the explicit formulae - we note that if $f$ has
polynomial q-expansions of deg $\leq n$, then $T_\ell(f)$ has expansions of
degree $\leq n\ell$.
\end{kaproof}

\begin{kathm}{Proposition}{1.11.3.}
Let $n\geq 2$ and $k\geq 2$, \uline{or} $3\leq n\leq 11$ and $k\geq 1$. For any
prime $\ell$ not dividing $n$, and \uline{any} $\Z[1/n]$-module $K$, there is a
necessarily unique endomorphism of the space of modular forms of weight $k$ and
level $n$, holomorphic at $\infty$, with coefficients in $K$, whose effect on
q-expansions is that given by the formulas (1.11.1.0-2).
\end{kathm}

\begin{kaproof}
By the base-changing theorem, we are reduced to the case $K = \Z[1/n]$. For a
modular form $f$ over $\Z[1/n]$, $T_\ell$ exists à priori over $\Z[1/n\ell]$,
but its q-expansions all have coefficients in $\Z[1/n,\zeta_n]$, so by \eqref{ka:1.6.2}
and \eqref{ka:1.9.1}, $T_\ell(f)$ is in fact a modular form over $\Z[1/n]$.
\end{kaproof}

\begin{kathm}{Corollary}{1.11.4.}
Let $k\geq 2$. For any prime $\ell$, and any $\Z$-module $K$, there is a
necessarily unique endomorphism of the space of modular forms of weight $k$ and
level one, holomorphic at $\infty$, whose effect on the q-expansion is that
given by the formulas (1.11.1.0-2).
\end{kathm}

\begin{kaproof}
Choose relatively prime integers $n,m\geq 3$, both prime to $\ell$, and view the
module of level one modular forms as the fibre-product of the diagram

\begin{equation}\tag{1.11.4.1}\label{ka:1.11.4.1}
\begin{tikzcd}
H^0(\Mbar_n,\w^{\otimes k}\tens(K\tens\Z[1/n])) \arrow[d] \\
H^0(\Mbar_{mn},(\w)^{\otimes k}\tens(K\tens\Z[1/nm])) & H^0(\Mbar_m,(\w)^{\otimes k}\tens(K\tens\Z[1/m])) \arrow[l]
\end{tikzcd}
\end{equation}

The desired $T_\ell$ is the fibre product of the $T_\ell$ constructed above on
this diagram.
\end{kaproof}

\sechead{1.12}{Applications to polynomial q-expansions; the strong q-expansion principle}

In this section we will admit the following result, a special case of
Swinnerton-Dyer's structure theorem (cf.\cite{serre73}, \cite{swinnerton-dyer-proceedings}), which will be proven later
(cf. \eqref{ka:4.4.1}).

\begin{kathm}{Result}{1.12.0.}
Let $n\geq 1$ be an integer, $K$ a field of characteristic $p\nmid n$, and $f$ a
modular form\index{modular form} over $K$ of level $n$ and weight $k\geq 1$, holomorphic at
infinity. Suppose $p-1\nmid k$. Then if all the q-expansions of $f$ at the cusps
of $\Mbar_n\tens K(\zeta_n)$ are \uline{constants}, $f = 0$.
\end{kathm}

Using this result, we will now prove

\begin{kathm}{Theorem}{1.12.1.}
Let $n$, $k\geq 1$ be integers, and suppose that $f$ is a modular form of level
$n$ and weight $k$, holomorphic at $\infty$, with coefficients in a
$\Z[1/n]$-module $K$. Suppose that for every prime $p$ such that $p-1|k$, the
endomorphism "multiplication by p" is injective on $K$. Then if all the
q-expansions of $f$ are polynomials in $q$, $f = 0$.
\end{kathm}

\begin{kaproof}
We begin by reducing to the case $n\geq 3$, using the diagram \eqref{ka:1.9.0.0} to
handle the case $n=1$, and the interpretation \hyperref[ka:1.9.0.1]{(1.9.1.1)} for $n=2$. We then
reduce to the case in which $n$ is divisible by $a = \prod_{p-1|k}p$; by
hypothesis $K\subset K[1/a]$, so we may \uline{replace} $K$ by $K[1/a]$ (using
the cohomology sequence \eqref{ka:1.6.2.2}), \uline{then} view $f$ as a modular form of
level $a\cdot n$ with coefficients in $K[1/a]$. Next we reduce to the case in
which $K$ is an artin local ring over $\Z[1/n]$, as explained in the proof of
\eqref{ka:1.6.1}. We will proceed by \uline{induction} on the least integer $b\geq 1$
such that $\mathfrak{m}^b = 0$, $\mathfrak{m}$ denoting the maximal ideal. Thus we begin with the
case in which $K$ is a \uline{field}.

Consider the finite-dimensional $K$-space $V$ of such modular forms, and choose
a basis $f_1,\ldots,f_r$ of $V$. Let $N$ be the maximum of the degrees of the
q-expansions of the $f_i$ at any of the cusps. At each cusp, record the
q-expansion of $F = \begin{pmatrix}f_1\\\vdots\\f_r\end{pmatrix}$:

\[
F(\mathrm{Tate}(q^n),\omega_{can},\alpha_n) = \sum_{i=0}^N A_i(\alpha_n)q^i, \qquad A_i = \begin{pmatrix}a_{i,1}(\alpha_n)\\\vdots\\a_{i,n}(\alpha_n)\end{pmatrix}
\]

Let $\ell$ be a prime number such that $\ell\nmid n$, $\ell > N$. Because $V$ is
stable under the Hecke operator\index{Hecke operator} $T_\ell$ (cf.\ref{ka:1.11}), we have a matrix equation
($C$ denoting an $r\times r$ matrix with coefficients in $K$),

\[
T_\ell(F) = C\cdot F.
\]

Passing to q-expansions gives the equation

\[
\sum_i(A_{\ell i}(\alpha_n'') + \ell^{k-1}A_{i/\ell}(\alpha_n'))q^i = C\cdot\sum_i A_i(\alpha_n)q^i
\]

whence, comparing coefficients of $q^{i\ell}$, we find the relation

\[
A_{\ell^2i}(\alpha_n'') + \ell^{k-1}A_{i\ell}(\alpha_n') = C\cdot A_{i\ell}(\alpha_n).
\]

But for $i\geq 1$, $i\ell > N$ and $i\ell^2 > N$, hence $A_{i\ell}(\alpha_n) = 0$
and $A_{\ell^2i}(\alpha_n'')=0$ (by definition of $N$). As $\ell$ is
invertible, we have $A_i(\alpha_n') = 0$ for each level n structure $\alpha_n$.
Hence each q-expansion of each $f_i$ is a constant, hence by \eqref{ka:1.12.0} each
$f_i = 0$. This concludes the proof in case $K$ is a field, and \uline{implies}
the case in which $K$ is a \uline{vector space} over a field, as vector spaces
have \uline{bases}.

Now consider the case of an Artin local ring K whose maximal ideal $\mathfrak{m}$
satisfies $\mathfrak{m}^{b+1} = 0$. By induction, $f$ becomes $0$ in $K/\mathfrak{m}^b$, hence by
the exact cohomology sequence \eqref{ka:1.6.2.2} associated to the exact sequence of
$\Z[1/n]$-modules $0\to\mathfrak{m}^b\to K\to K/\mathfrak{m}^b\to 0$, $f$ \uline{comes from} a
form with coefficients in $\mathfrak{m}^b$. But as $\mathfrak{m}^{b+1} = 0$, $\mathfrak{m}^b$ is a
(finite-dimensional!) vector space over the residue field $K/\mathfrak{m}$, and the
previous case of a field applies.
\end{kaproof}

\begin{kathm}{Corollary}{1.12.2.}
(Strong q-expansion principle\index{q-expansion principle}) Let $n$, $k\geq 1$, and let
$a = \prod_{p-1|k}p$. Let $K$ be a $\Z[1/an]$-module of which $L\subset K$ is a
$\Z[1/an]$-submodule, and $f$ a modular form of level $n$ and weight $k$,
holomorphic at $\infty$, such that at \uline{each} cusp, all but finitely many
of its q-expansion coefficients lie in
$L\tens_{\Z[1/n]}\Z[1/n,\zeta_n]$. Then $f$ \uline{is} a modular form with
coefficients in $L$.
\end{kathm}

\begin{kaproof}
Apply the theorem to the image of $f$ as modular form with coefficients in
$K/L$.
\end{kaproof}

\sechead{1.13}{Review of the modular scheme associated to $\Gamma_0(p)$}

For each integer $n\geq 3$ prime to $p$, the functor "isomorphism classes of
elliptic curves\index{elliptic curve} with level n structure and a finite flat subgroup (scheme) of
rank p" is \uline{representable}, by a scheme $M_{n,p}$, which is an affine
curve over $\Z[1/n]$; it is a \uline{regular} scheme, but it \uline{fails} to
be smooth over $\Z[1/n]$ precisely at the finitely closed points on $M_{n,p}$
corresponding to \uline{super-singular} elliptic curves in characteristic $p$.
The projection "forget the subgroup of rank p" makes $M_{n,p}$ finite and flat
over $M_n$ of degree $p+1$.

We define $\Mbar_{n,p}$ to be the normalization of $\Mbar_n$ in $M_{n,p}$; it is
a regular scheme, proper and flat over $\Z[1/n]$. The difference
$\Mbar_{n,p} - M_{n,p}$ is finite and étale over $\Z[1/n]$, and over
$\Z[1/n,\zeta_n]$ it is a disjoint union of sections, called the cusps of
$\Mbar_{n,p}$, two of which lie over each cusp of $\Mbar_n$, and exactly one of
which is étale over $\Mbar_n$.

The completion of $\Mbar_{n,p}\tens\Z[1/n,\zeta_n]$ along any of the cusps is
isomorphic to $\Z[1/n,\zeta_n][[q]]$. The universal elliptic curve with level n
structure and subgroup of order p over $\Z[1/n,\zeta_n]((q))$, viewed as a
punctured disc around an unramified cusp, is the Tate curve\index{Tate curve}
$\mathrm{Tate}(q^n)$ with the level n structure corresponding to the underlying
cusp of $\Mbar_n$, and the subgroup $\mu_p$. Over one of the ramified cusps, the
inverse image is the Tate curve $(q^{np})$, with the induced ($q\mapsto q^p$)
level n structure from the cusp of $\Mbar_n$ below, and with the subgroup
generated by $q^n$.

The automorphism of $M_{n,p}$ given by
$(E,\alpha_n,H)\mapsto(E/H,\pi(\alpha_n),{}_pE/H)$ ($\pi\colon E\to E/H$
denoting the projection, and $\pi(\alpha_n)$ the level n structure explained in
\eqref{ka:1.11.0.0}) extends to an automorphism of $\Mbar_{n,p}$ which interchanges the
two sorts of cusps.
